% =============================================================================
%  导言区：中译注本《三维镜像对称》
%  原文：Ben Webster, Philsang Yoo,
%        "3-dimensional mirror symmetry", arXiv:2308.06191v1 [math.RT],
%        Notices of the AMS（本译本据 arXiv 版，正文 23 页）
%  编译：xelatex main.tex（三遍）
% =============================================================================
\documentclass[11pt,a4paper]{article}

\usepackage[UTF8, fontset=fandol, heading=true]{ctex}
\usepackage{amsmath}
\usepackage{amssymb}
\usepackage{amsthm}
\usepackage{mathtools}
\usepackage{bm}
\usepackage[a4paper,left=2.5cm,right=2.5cm,top=2.6cm,bottom=2.6cm]{geometry}
\usepackage{graphicx}
\usepackage{xcolor}
\usepackage[most]{tcolorbox}
\usepackage{enumitem}
\usepackage{booktabs}
\usepackage{longtable}
\usepackage{array}
\usepackage{multirow}
\usepackage{tikz}
% 原文 main.tex:12–17 的图样式，逐字照搬
\tikzset{
  MyPersp/.style={scale=1.8,x={(-0.8cm,0cm)},y={(0cm,.3cm)},z={(0cm,1cm)}},
  MyPoints/.style={fill=white,draw=black,thick}
}
\usepackage{tikz-cd}
\usepackage{tqft}
\usepackage{wrapfig}
\usepackage[font=small,labelfont=bf,labelsep=quad]{caption}
\usepackage[style=alphabetic,url=false,isbn=false,doi=false,maxbibnames=99]{biblatex}
% 原包只带 .bbl（3.2 格式）与 2026 年的 biblatex（3.3）不兼容，故由 .bbl 反推 main.bib，
% 用 biber 重建文献表；见 BRIEF.md「文献表」条。
\addbibresource{main.bib}
\AtEveryBibitem{%
  \clearfield{urlyear}%
  \clearfield{urlmonth}%
}
\usepackage[colorlinks=true,
            linkcolor=blue!45!black,
            citecolor=blue!45!black,
            urlcolor=blue!45!black]{hyperref}

\linespread{1.15}
\setlength{\emergencystretch}{3em}
\allowdisplaybreaks[3]
\pagestyle{plain}
\renewcommand{\figurename}{图}
\renewcommand{\refname}{参考文献}
\renewcommand{\contentsname}{目录}
\renewcommand{\proofname}{证明}
\renewcommand{\abstractname}{摘要}
\renewcommand{\tablename}{表}

% 原文以 \emph / {\bf } 标出强调；中文以加粗对应
\renewcommand{\emph}[1]{{\bfseries #1}}

% ---------------------------------------------------------------------------
% 原文的自定义宏（与原文同名，见原文 preamble；译手只允许使用这些）
% ---------------------------------------------------------------------------
\newcommand{\Sec}{\operatorname{Sec}}
\newcommand{\C}{\mathbb{C}}
\newcommand{\R}{\mathbb{R}}
\newcommand{\Z}{\mathbb{Z}}
\newcommand{\Spec}{\operatorname{Spec}}
\newcommand{\Hom}{\operatorname{Hom}}
\newcommand{\Bv}{\mathbf{v}}
\newcommand{\Bw}{\mathbf{w}}
\newcommand{\Sym}{\operatorname{Sym}}
\newcommand{\obs}{\operatorname{Obs}}
\newcommand{\NO}{N_{\mathcal{O}}}
\newcommand{\NK}{N_{\mathcal{K}}}
\newcommand{\GO}{G_{\mathcal{O}}}
\newcommand{\GK}{G_{\mathcal{K}}}
\newcommand{\XO}{X_{\mathcal{O}}}
\newcommand{\XK}{X_{\mathcal{K}}}
\newcommand{\XA}{X_n^{(\mathrm{A})}}
\newcommand{\XB}{X_n^{(\mathrm{B})}}
\newcommand{\XS}{X_n^{(*)}}
\newcommand{\fM}{\mathfrak{M}}
\newcommand{\disk}{\mathbb{D}}
\newcommand{\pdisk}{\mathbb{D}^{\circ}}
\newcommand{\rkone}{M_{n\times n}^{\mathrm{rk}\leq 1}(\C)}
\newcommand{\Symp}{\mathrm{Symp}}
\newcommand{\Ham}{\mathrm{Ham}}
\newcommand{\KR}{\mathrm{K}}
\newcommand{\id}{\mathrm{id}}
% 原文 orig/main.tex:65 / :84 定义、本库首稿漏收的两个宏（2026-11-22 补）：
\newcommand{\alg}{\mathcal{A}}
\newcommand{\sPoin}{\mathfrak{a}}

% ---------------------------------------------------------------------------
% 定理类环境：与原文编排一致（每节内一个共享计数器），
% 编号如「定理 3.1」「定义 3.2」，与原文逐条对应。
% ---------------------------------------------------------------------------
\newtheoremstyle{cnth}%
  {6pt}% above
  {6pt}% below
  {\normalfont}% body
  {}% indent
  {\bfseries}% head font
  {.}% punctuation
  {0.5em}% space after head
  {}% head spec

\theoremstyle{cnth}
\newtheorem{theorem}{定理}[section]
\newtheorem{proposition}[theorem]{命题}
\newtheorem{corollary}[theorem]{推论}
\newtheorem{lemma}[theorem]{引理}
\newtheorem{assumption}[theorem]{假设}
\newtheorem{definition}[theorem]{定义}
\newtheorem{ex}[theorem]{例}
\newtheorem{note}[theorem]{注}
\newtheorem{ques}[theorem]{问题}
\newtheorem{conjecture}[theorem]{猜想}
\newtheorem{physics}[theorem]{物理动机}

\renewcommand{\qedsymbol}{$\square$}

% ---------------------------------------------------------------------------
% 译注方框：橙色、编号；可用 \showyizhufalse 关闭
% ---------------------------------------------------------------------------
\newif\ifshowyizhu
\showyizhutrue
\newcounter{yizhu}

\newenvironment{yizhu}[1][]{%
    \refstepcounter{yizhu}%
    \par\medskip
    \ifshowyizhu
        \begin{tcolorbox}[enhanced, breakable,
            colback=orange!4!white, colframe=orange!55!black,
            boxrule=0.5pt, arc=2pt, left=6pt, right=6pt, top=4pt, bottom=4pt,
            title={\small\textbf{译注 \theyizhu\if\relax\detokenize{#1}\relax\else： #1\fi}},
            fonttitle=\bfseries, coltitle=black, colbacktitle=orange!12!white,
            toptitle=1pt, bottomtitle=1pt]
    \fi
}{%
    \ifshowyizhu
        \end{tcolorbox}%
    \fi
    \par\smallskip
}

% 行内译者小注
\newcommand{\zy}[1]{\footnote{译者注：#1}}

% 术语表条目：中文（English）
\newcommand{\term}[2]{#1（#2）}
